\input{../article_base.tex}
\title{תורת ההסתברות 2 --- סיכום}
\setcounter{secnumdepth}{2}
% chktex-file 9
% chktex-file 17

\usepackage{fancyhdr}
\pagestyle{fancy}
\renewcommand{\headrulewidth}{0pt}

\begin{document}
\maketitle
\maketitleprint[teal]{}

\tableofcontents

\section{שיעור 1 --- 25.3.2026}

\subsection{מבוא הקורס}
הקורס מבוסס על הספר Probability with Martingales של Williams.
במהלך הקורס נלמד שלושה נושאים, קצת סותרים באיזשהו מובן.
נתחיל מללמוד הסתברות בהתבסס על תורת המידה, כך נבסס תוצאות שנלמדו בקורס הקודם.
בחלק השני נלמד נושאים מתקדמים בהסתברות, כדוגמת משפט הגבול המרכזי (אותו ראינו אבל לא הוכחנו), פורייה מרטינגלים ושרשרות מרקוב.
בחלק האחרון נגיע לנושאים ומחקר בעולם ההסתברות, נלמד על הילוכים מקריים, פרוקלציה וכן הלאה.

נעבור לדוגמות דווקא מהחלק האחרון, כדי להדגים רעיונות אותם אנו רוצים לראות.
\begin{example}[הילוכים מקריים]
	באנגלית (Simple) Random Walk.
	יהיו $X_1, X_2, \ldots$ משתנים מקריים שווי־התפלגות כך ש־$X_n \sim U(\{-1, 1\})$ ונגדיר $Y_n = \sum_{i = 1}^n X_i$.
	הסיבה שאנחנו קוראים למושג הזה הילוך היא שהערכים המתקבלים מהסכימה מהלכים על הישר הממשי, אנו מעוניינים לשאול שאלות כמו $\PP(Y_n = 0) = 0$.
	אנו גם יכולים לבדוק את הערך $\PP(\exists n > 0\ Y_n = 0)$, ברור כי ערך זה הוא לפחות חצי שכן הוא מקיים $\ge \PP(Y_2 = 0) = \frac{1}{2}$.
	למעשה ערך זה שווה ל־$\lim_{m \to \infty} \PP(\exists 0 < n < m\ Y_n = 0) = 1$, ונאמר שבמקרה זה ההילוך המקרי נשנה.
	אם $X_n$ היו מתפלגים אחרת, לדוגמה $\PP(X_n = 1) = p$ עבור $0 \le p \le 1$ אז $\PP(\exists n > 0\ Y_n = 0) < 1$.

	נבחן את $\PP(Y_{2n} = 0)$ זוהי ההסתברות שיש בדיוק $n$ משתנים $X_k$ המקבלים $1$,
	\begin{align*}
		\PP(Y_{2n} = 0)
		& = \PP(\exists I \subseteq [2n], |I| = n\ \forall i\ X_i = 1 \iff i \in I) \\
		& = \sum_{|I| = n, I \subseteq [2n]} \PP(\forall i\ X_i = 1 \iff i \in I) \\
		& = \sum_{|I| = n, I \subseteq [2n]} \frac{1}{2^{2n}} \\
		& = \binom{2n}{n} \cdot \frac{1}{2^{2n}} \\
		& = \frac{(2n)!}{{(n!)}^2 2^{2n}} \\
		& \overset{(1)}{\approx}  \frac{\sqrt{2 \pi \cdot 2n} {(\frac{2n}{e})}^{2n}}{{(\sqrt{2 \pi n} {(\frac{n}{e})}^n)}^2 \cdot 2^{2n}} \\
		& = \frac{\sqrt{4 \pi n}}{2 \pi n}
		= \frac{1}{\sqrt{\pi n}}
	\end{align*}
	כאשר $(1)$ נובע מנוסחת סטירלינג, כלומר ש־$k! \approx \sqrt{2 \pi k} {(\frac{k}{e})}^k$.
	לכן $\sum_{n = 1}^\infty \PP(Y_{2n} = 0) = \infty$.
	קיבלנו את התוצאה שמספר החזרות לראשית הוא אינסוף.

	נעבור למקרה $\PP(X_i = 1) = p > \frac{1}{2}$.
	אי־שוויון הופדינג טוען שאם $Z_1, \ldots, Z_n$ משתנים מקריים בלתי־תלויים עם $\EE(Z_i) = 0$ לכל $i$ וכן ש־$|Z_i| \le 1$, אז לכל $a > 0$ מתקיים,
	\[
		\PP(\sum_{i = 1}^n Z_i \ge a)
		\le \exp(-\frac{a^2}{2n})
	\]
	צ'בישב היה נותן לנו חסם $\le \frac{n}{a^2}$ במקום.
	נגדיר $Z_i = \frac{X_i - (2p - 1)}{2}$ כדי לנרמל ונקבל שמתקיים $\sum_{i = 1}^k Z_i = \frac{1}{2} \sum_{i = 1}^k X_i - n(p - \frac{1}{2})$, אז מתקיים,
	\begin{align*}
		\PP(Y_n = 0)
		& = \PP(\sum_{i = 1}^n X_i = 0) \\
		& = \PP(\sum_{i = 1}^n Z_1 = -n (p - \frac{1}{2})) \\
		& \le \PP(\sum_{i = 1}^n -Z_i \ge n(p - \frac{1}{2})) \\
		& \le \exp(-\frac{n^2 {(p - \frac{1}{2})}^2}{2n}) \\
		& = \exp(-\frac{{(p - \frac{1}{2})}^2}{2} \cdot n)
	\end{align*}
	כלומר ההסתברות דועכת אקספוננציאלית, ובהתאם $\sum_{n = 1}^\infty \PP(Y_n = 0) < \infty$.
	מצאנו שיש הבדל מהותי בין המקרים $p = \frac{1}{2}$ והמקרים $p \ne \frac{1}{2}$.

	אם נסמן $q = \PP(\forall n > 0\ Y_n = 0)$, כלומר ההסתברות שלעולם לא נחזור לאפס.
	התפלגות מספר החזרות לראשית היא גאומטרית $\operatorname{Geo}(q)$.
	אם $q = 0$ אז בהסתברות $1$ חוזרים לראשית אינסוף פעמים.
	אם $q > 0$ אז תוחלת מספר החזרות לראשית היא $\frac{1}{q} < \infty$.
	כש־$p = \frac{1}{2}$ התוחלת היא אינסוף ואחרת התוחלת סופית.
	לכן אם $p = \frac{1}{2}$ אז $q = 0$ (נשנה) ואם $p \ne \frac{1}{2}$ אז $q > 0$ (חולף).
	זוהי לא הוכחה פורמלית, אבל היא מבהירה לנו את הקונספט הכללי עליו אנו מדברים.
\end{example}
\begin{definition}[הילוך מקרי פשוט על גרף]
	באנגלית Simple Random Walk (SRW) על גרף.
	נניח ש־$G = (V, E)$ גרף פשוט וקשיר עם דרגות סופיות ויהי $v_0$ קודקוד כלשהו.
	SRW היא סדרת קודקודים מקריים בגרף, כאשר $X_0 = v_0$ כמעט תמיד, בהינתן ש־$X_0 = v_0, X_1 = v_1, \ldots, X_n = v_n$ מגרילים באופן אחיד בין השכנים של $v_n$.
\end{definition}
\begin{example}
	עתה נוכל לדבר על $q = \PP(\forall n > 0\ X_n \ne v_0)$, כלומר הפעם ההילוך הוא לאורך גרף, אם $q = 0$ אז ההילוך נשנה ואם $q > 0$ אז ההילוך חולף.

	אם $G = \ZZ^2$ אז ההילוך נשנה, כדי להראות זאת עלינו להראות שתוחלת מספר החזרות לראשית הוא אינסופי, נבחין כי $\sum_{n = 1}^\infty \PP(X_n = v_0) = \infty$ שכן $\PP(X_n = v_0) \approx \frac{c}{n}$.
	כדי לפתור את הבעיה הזו אנו נבחן את המקרה שבו ההילוך הוא בלתי תלוי על שני הצירים, התוצאה תהיה תנועה על רשת מסוימת שנראית בדיוק כמו $\ZZ^2$, ולכן נקבל את המבוקש כמכפלה של בלתי־תלויים.

	מה יקרה במקרה $G = \ZZ^d$ עבור $d > 2$? השיטה של $d = 2$ לא עובדת יותר, אבל הטענה עדיין נכונה, כלומר $\PP(X_n = v_0) \approx \frac{c}{n^\frac{d}{2}}$ ולכן $\sum_{n = 1}^\infty \PP(X_n = v_0) = \infty$ ולכן ההילוך חולף.
\end{example}
מה שראינו עתה הוא תוצאה בסיסית אחת בעולם שלם.
באופן כללי האם גרף הוא נשנה או לא היא תכונה של הגרף ולא של נקודת ההתחלה.

\subsection{פרקולציה}
נעבור לנושא אחר אך קשור.
נניח שיש לנו גרף $G$ קשיר ואינסופי, לדוגמה $\ZZ^2$.
לכל קשת נגריל בהסתברות $p$ צביעה בכחול או אדום, באופן בלתי תלוי.
יש לנו כמה שאלות שאנחנו יכולים לשאול, לדוגמה מה ההסתברות שהגרף הכחול קשיר, הסתברות זו היא $0$ אם $p < 1$, זאת שכן ההסתברות שנקודה כלשהי נתונה היא מבודדת היא $1 - p^4$.
נעבור לשאלה מה ההסתברות שבגרף הכחול יש רכיב קשירות אינסופי, זוהי שאלה מורכבת הרבה יותר, נסמן $\theta(p)$ כהסתברות ביחס ל־$p$, ונבחין ש־$\theta$ היא מונוטונית עולה חלש.
יהיו $0 \le p_1 < p_2 \le 1$ ונצבע כל קשת בהסתברות $p_1$ בכחול כהה, בהסתברות $p_2 - p_1$ נצבע בכחול בהיר וב־$1 - p_2$ לא נצבע בכלל.
אם נבחן את $\theta(p_1)$ נקבל את המצב המקורי עם $p_1$, אם נבחן את שני גווני הכחול נקבל את $\theta(p_2)$, נגדיר את המורע $A$ להיות קיום רכיב קשירות אינסופי כחול בהיר ו־$B$ רכיב קשירות כזה כחול (שני הגוונים).
אז $A \subseteq B$ ולכן $\theta(p_1) = \PP(A) \le \PP(B) = \theta(p_2)$.
נקרא לרעיון זה צימוד (Coupling).
משפט קולומורוב אומר שאם $A$ מאורע זנב (לא מושפע משום קבוצה סופית של קורדינטות) אז $\PP(A) \in \{0, 1\}$, נסיק ש־$\theta(p) \in \{0, 1\}$ ומונוטונית ונסמן את נקודת השינוי ב־$p_c$.
האם $0 < p_c < 1$? מה ערך $\theta(p_c)$? אלו הן שאלות שאנו עדיין לא יכולים לענות עליהן עכשיו.
אם $G = \ZZ$ אז $p_c = 1$.
התשובה היא ש־$p_c = \frac{1}{2}, \theta(\frac{1}{2}) = 0$ עבור $G = \ZZ^2$ וזה מקרה נדיר שבו יש לנו ערך ולא רק חסם כללי.

\section{שיעור 2 --- 15.4.2026}
\subsection{ממידה להסתברות}
עתה כשאנו מכירים את תורת המידה, אנו יכולים לבסס מחדש את תורת ההסתברות עם כלים של מידה, לדוגמה מרחב הסתברות הוא מרחב מידה כך ש־$\mu(\Omega) = 1$.
בנוסף נגדיר שמשתנה מקרי $X : \Omega \to \RR$ כפונקציה מדידה.
נאמר גם שנעסוק רק בקבוצות מדידות בורל.
גם הגדרנו את $\PP_X(S) = \PP(X^{-1}(S))$ להיות התפלגות $X$, וכן $\EE(X) = \int_{\Omega} X\ d\PP = \int_{\RR} 1\ d \PP_X$.
עתה ניגש למלאכה הפורמלית.

נניח מעתה ש־$\Omega$ קבוצה לא ריקה
\begin{definition}[אלגברה]
	$\Ff \subseteq \Pp(\Omega)$ היא אלגברה על $\Omega$ אם מתקיים,
	\begin{enumerate}
		\item $\emptyset \in \Ff$
		\item לכל $A \in \Ff$ גם $A^C \in \Ff$
		\item $\Ff$ סגורה לאיחוד זר סופי
	\end{enumerate}
	נאמר ש־$\Ff$ היא $\sigma$־אלגברה אם היא סגורה לאיחוד זר בן־מניה.
\end{definition}
\begin{definition}[מרחב מדיד]
	מרחב מדיד $(\Omega, \Ff)$ הוא קבוצה $\Omega$ ו־$\sigma$־אלגברה עליה.
\end{definition}
\begin{example}
	איחודים סופיים של קטעים $(a, b) \subseteq (0, 1]$ היא אלגברה.
\end{example}
\begin{example}
	עבור $\Omega = {\{0, 1\}}^\NN$ איחודים סופיים של קבוצות צילינדריות, כלומר,
	\[
		\forall a \in {\{ 0, 1 \}}^\NN,
		\quad
		C_a
		= \{ x \in \Omega \mid \forall 1 \le i \le n,\ x = a_i \}
	\]
	היא אלגברה.
	במקרה זה לדוגמה $C_{0, 1, 0}$ היא קבוצת כל הסדרות שמתחילות ברישא $(0, 1, 0)$.
\end{example}
\begin{definition}[סיגמא־אלגברה נוצרת]
	לכל $S \subseteq \Pp(\Omega)$ נגדיר את $\sigma(S)$ להיות ה־$\sigma$־אלגברה המינימלית שמכילה את $S$.
\end{definition}
נבחין כי האופן שקול אפשר להגדיר את ה־$\sigma$־אלגברה הנוצרת באופן קונסטרוקטיבי.
\begin{example}
	עבור $S = \{ A \}$ מתקיים $\sigma(S) = \{ \emptyset, A, A^C, \Omega \}$.
\end{example}
\begin{example}
	עבור $A \subseteq \NN$ נאמר שיש לה צפיפות $a \in [0, 1]$ אם,
	\[
		\lim_{n \to \infty} \frac{|A \cap [n]|}{n}
		= a
	\]
	נגדיר את $\Ff$ כאוסף כל הקבוצות בעלות צפיפות, האם $\Ff$ היא $\sigma$־אלגברה?
	התשובה היא שלא שכן כל יחידון הוא בעל צפיפות.
	האם $\Ff$ היא אלגברה?
\end{example}
\begin{definition}[מידת הסתברות]
	בהינתן מרחב מדיד $(\Omega, \Ff)$, כלומר $\Ff$ היא $\sigma$־אלגברה,
	פונקציה $\PP : \Ff \to [0, 1]$ היא מידת הסתברות אם,
	\begin{enumerate}
		\item $\PP(\emptyset) = 0$
		\item $\PP(A) \ge 0$ לכל $A \in \Ff$
		\item $\sigma$־אדיטיביות: אם ${\{ A_n \}}_{n = 1}^{\infty} \subseteq \Ff$ זרות רז $\PP(\bigcup_{n \in \NN} A_n) = \sum_{n = 1}^\infty \PP(A_n)$
		\item $\PP(\Omega) = 1$
	\end{enumerate}
	אם $\mu$ מוגדרת על אלגברה $\Ff$ ומקיימת סגירות לאיחוד זר בן־מניה אז נקרא לה \textbf{קדם־מידה}.
\end{definition}
\begin{theorem}[ההרחבה של קרתאודורי]
	אם $\Ff_0$ אלגברה על $\Omega$ ו־$\mu_0$ קדם־מידה סופית אז קיימת מידה יחידה וסופית $\mu$ על $\Ff = \sigma(\Ff_0)$ שמרחיבה את $\mu_0$.
\end{theorem}
\begin{example}
	עבור $\Omega = {\{0, 1\}}^\NN$ ו־$\Ff_0$ האיחודים הסופיים של צילינדרים, $\Ff = \sigma(\Ff_0)$ נגדיר,
	\[
		\forall a \in \Omega^{< \omega},
		\quad
		\mu(C_a)
		= 2^{-|a|}
	\]
	כלומר נקבל את האורך של הסדרה הסופית $a$, ונקבל את ההסתברות של הטלת מטבע מספר סופי של פעמים וקבלת הסדרה $a$.
	המשפט נותן לו מידת הסתברות $\mu$ המרחיבה את $\mu_0$, עתה נוכל לדבר על מאורעות ממידה אפס. \\
	מה הקשר בין $\mu$ למידת לבג על $(0, 1]$?
	התשובה היא שהמידה של סדרה שווה למידה שלו כמספר בייצוג בינארי עשרוני במידת לבג, ואלו מרחבים שקולים עד כדי מידה 0.
\end{example}
\begin{definition}[גבול עליון ותחתון של מאורעות]
	עבור ${\{ A_n \}}_{n = 1}^\infty \subseteq \Ff$ סדרת מאורעות (מדידות) נגדיר,
	\[
		\limsup A_n = \bigcap_{m = 1}^{\infty} \bigcup_{n = m}^{\infty} A_n,
		\qquad
		\liminf A_n = \bigcup_{m = 1}^{\infty} \bigcap_{n = m}^{\infty} A_n
	\]
\end{definition}
ההצדקה להגדרה זו היא שאם נגדיר את סדרת הפונקציות $\indicator_{A_n}$ אז נקבל,
\[
	\limsup_{n \to \infty} \indicator_{A_n}
	= \indicator_{\limsup_{n \to \infty} A_n}
\]
ניזכר במספר טענות,
\begin{lemma}[הראשונה של בורל־קנטלי]
	אם ${\{ A_n \}}_{n = 1}^\infty$ כך ש־$\sum_{n \in \NN} < \infty$ אז $\PP(\lim_{n \to \infty} A_n) = 0$.
\end{lemma}
נראה דוגמה נגדית לגרירה ההפוכה,
\begin{example}
	אם נבחר את $A_n = [0, \frac{1}{n}]$ אז נקבל $\PP(A_n) = \frac{1}{n}$ ולכן הסכום מתבדר אבל גם $\limsup_{n \to \infty} A_n = \{ 0 \}$ ומהסתברות 0.
\end{example}
\begin{lemma}
	אם ${\{ A_n \}}_{n = 1}^\infty \subseteq \Ff$ מאורעות בלתי־תלויים כך ש־$\sum_{n \in \NN} \PP(A_n) = \infty$ אז $\PP(\limsup_{n \to \infty} A_n) = 1$.
\end{lemma}
בקורס הקודם הגדרנו אי־תלות בדרך ישירה,
\begin{definition}[אי־תלות של מאורעות]
	אם ${\{ A_n \}}_{n \in \NN} \subseteq \Ff$ סדרת מאורעות, אז נאמר שהיא בלתי־תלויה אם לכל $I \subseteq \NN$ סופית מתקיים,
	\[
		\PP(\bigcap_{i \in I} A_i)
		= \prod_{i \in I} \PP(A_i)
	\]
\end{definition}
אבל נוכל להרחיב את ההגדרה לאי־תלות כללית יותר,
\begin{definition}[אי־תלות סיגמא־אלגברה]
	אם $(\Omega, \Ff, \PP)$ מרחב הסתברות ו־$\Ff_1, \Ff_2$ הן תת־$\sigma$־אלגברות,
	אז נאמר שהן בלתי־תלויות אם לכל $A_1 \in \Ff_1, A_2 \in \Ff_2$ מתקיים,
	\[
		\PP(A_1 \cap A_2)
		= \PP(A_1) \cdot \PP(A_2)
	\]
	מאורעות $A_1, A_2$ יקראו בלתי־תלויים אם $\sigma(A_1)$ ו־$\sigma(A_2)$ הן בלתי־תלויות.
\end{definition}
\begin{definition}[אי־תלות קבוצת סיגמא־אלגברות]
	אם $I$ קבוצת אינדקסים כלשהי, ${\{ \Ff_i \}}_{i \in I}$ תת־$\sigma$־אלגברות במרחב,
	אז נאמר שהן בלתי־תלויות אם לכל $J \subseteq I$ סופית מתקיים ולכל בחירה של $A_j \in \Ff_j$ לכל $j \in J$ מתקיים,
	\[
		\PP(\bigcap_{j \in J} A_j)
		= \prod_{j \in J} \PP(A_j)
	\]
	מאורעות ${\{ A_i \}}_{i \in I}$ יקראו בלתי־תלויים אם ${\{ \sigma(A_i) \}}_{i \in I}$ בלתי־תלויה.
\end{definition}
נחזור ללמה השנייה ונוכיח אותה.
\begin{proof}
	נבחין כי מתקיים,
	\[
		\limsup_{n \to \infty} A_n
		= {\left(\liminf_{n \to \infty} A_n^C\right)}^C
	\]
	ולכן אנו רוצים להראות שמתקיים $\PP(\liminf_{n \to \infty} A_n^C) = 0$, כלומר $\PP(\bigcup_{m = 1}^{\infty} \bigcap_{n = m}^{\infty} A_n^C) = 0$, אבל מרציפות,
	\[
		\PP(\bigcup_{m = 1}^{\infty} \bigcap_{n = m}^{\infty} A_n^C)
		\le \sum_{m = 1}^\infty \PP(\bigcap_{n = m}^\infty A_n^C)
	\]
	אבל עבור סכומים חלקיים נקבל,
	\[
		\lim_{m \to \infty} \PP(\bigcap_{n = m}^M A_n^C)
		= \lim_{m \to \infty} \prod_{n = m}^M (1 - \PP(A_n))
		\le \lim_{n \to \infty} \exp\left(- \sum_{n = m}^M \PP(A_n)\right)
		\to 0
	\]
	כאשר אי־השוויון האחרון נובע מאי־שוויון ברנולי $1 + x \le e^x$.
	בהתאם קיבלנו את המבוקש.
\end{proof}
\begin{definition}[סיגמא־אלגברת זנב]
	נניח ש־$(\Omega, \Ff)$ מרחב מדיד ו־${\{ \Ff_n \}}_{n \in \NN}$ סדרת תת־$\sigma$־אלגברות.
	נגדיר את $\sigma$־אלגברת הזנב של $\{ \Ff_n \}$ על־ידי,
	\[
		\Tt
		= \bigcap_{m = 1}^\infty \sigma(\Ff_m, \Ff_{m + 1}, \ldots)
	\]
	כאשר נגדיר $\sigma(A_1, A_2, \ldots) = \sigma(A_1 \cup A_2 \cup \cdots)$.
	למאורע $A \in \Tt$ נקרא מאורע זנב.
\end{definition}
\begin{remark}
	ההגדרה של $\sigma$־אלגברת זנב לא תלויה בסדר $\{ \Ff_n \}$.
\end{remark}
\begin{example}
	נגדיר שוב $\Omega = {\{ 0, 1 \}}^\NN$ וכן $\Ff_n = \sigma(B_n)$ כאשר $B_n = \{ x \in \Omega \mid x_n = 0 \}$.
	נבחין כי $A  = \{ (0, 0, \ldots) \} \notin \Tt$, שכן אם $B \in \Tt_2$ ו־$y \in B$ אז גם $y' \in B$ כאשר $y = y'$ מלבד הקורדינטה הראשונה.
	לעומת זאת, המאורע $B = \{ x \in \Omega \mid \exists m \forall n > m,\ x_n = 0 \}$ מקיים $B \in \Tt_m$.
\end{example}
\begin{theorem}[חוק 0,1 של קולמוגורוב]\label{theorem_cologomorov_law}
	אם $(\Omega, \Ff, \PP)$ מרחב הסתברות ו־$\Ff_n$ סדרת תת־$\sigma$־אלגברות בלתי־תלויה, ו־$A \in \Tt$ מאורע זנב, אז,
	\[
		\PP(A) \in \{0, 1\}
	\]
\end{theorem}
\begin{example}
	נבחן את הגרף $\ZZ^2$ עם צביעה אקראית בשני צבעים, ושאלנו בשיעור הקודם מה ההסתברות שיש רכיב קשירות אינסופי באחד הצבעים,
	עתה נטען שזהו מאורע זנב.
\end{example}

\section{שיעור 3 --- 29.4.2026}
נמשיך לעסוק במשפט\ \ref{theorem_cologomorov_law} ונראה דוגמה למקרה שהמשפט מייצג.
\begin{example}
	יהי $G = (V, E)$ גרף אינסופי קשיר עם דרגות סופיות, ונגדיר $\Omega = {\{ 0, 1 \}}^E$ כמרחב מכפלה של צביעות, כאשר ההסתברות לצביעה $1$ היא $p$ ו־$1 - p$ עבור $0$.
	עבור $w \in \Omega$ נגדיר $E' = \{ e \in E \mid w_e = 1 \}$ ונגדיר את תת־הגרף $(V, E')$.
	נסמן $F_e = \sigma(A_e)$ עבור $A_e = \{ w \in \Omega \mid w_e = 1 \}$.
	קיים ב־$(V, E')$ רכיב קשירות אינסופי הוא מאורע זנב.
	נבחין כי $A \in T_n$ אם ורק אם $x, y$ שונים רק ב־$n - 1$ הקורדינטות הראשונות אז או ש־$x, y \in A$ או ש־$x, y \notin A$ (תרגיל).
	נעבור לדון במקרה של $\ZZ^2$.
	נסיק שאם $f(p) = \cap_p(B) \cap \{0, 1\}$ אז $f$ פונקציה עולה חלש וכן $f(0) = 0, f(1) = 1$ ולכן קיים $0 \le p_c \le 1$ נקודת אי־רציפות.
	אם $B^C$ כל המאורעות עם רכיבי קשירות סופיים בלבד ו־$D_v$ המאורע ש־$v$ נמצא ברכיב קשירות סופי אז נקבל ש־$B^C = \bigcap_{v \in V} D_v$ וכן $D_v = \biguplus_{i \in \NN} D_v^i$.
	נטען ש־$\frac{1}{3} \le p_c$ אם ורק אם לכל $p < \frac{1}{3}$ מתקיים $f(p) = 0$.
	יהי $p < \frac{1}{3}$, אז הסתברות המאורע שקיים מסלול אינסופי פשוט שמתחיל ב־$v$ ונרצה למצוא את ההסתברות שקיים מסלול כזה באורך $n$.
\end{example}
\begin{proposition}
	תהי $\Ff$ $\sigma$־אלגברה ו־$\Ff_0$ אלגברה כך ש־$\Ff = \sigma(\Ff_0)$.
	לכל $A \in \Ff$ ולכל $\varepsilon > 0$ קיימת $A_0 \in \Ff_0$ כך ש־$0 \le \PP(A \triangle A_0) < \varepsilon$.
\end{proposition}
\begin{exercise}
	הוכיחו טענה זו.
\end{exercise}
יהי $A \in T$ ויהי $\varepsilon > 0$ אז,
\[
	A \in \sigma(\bigcup_{m \in \NN} \sigma(\bigcup_{k = 1}^m \Ff_k))
\]
לפי הטענה קיימת $A_0 \in \bigcup_{m \in \NN} \sigma(\cdots)$ כך ש־$\PP(A \triangle A_0) < \varepsilon$.
קיים $m_0$ כך ש־$A_0 \in \sigma(\bigcup_{k = 1}^{m_0} \Ff_k)$ ו־$A \in T_{m_0}$ בהתאם.
לכן $\PP(A \cap A_0) = \PP(A) \PP(A_0)$, נסמן $p = \PP(A)$ ולכן $\forall \varepsilon > 0, p - \varepsilon \le p(p + \varepsilon) = p^2 + p \varepsilon$.

\begin{theorem}[אינווריאנטיות להזזה בהסתברות]
	תהי $\Omega = {\{ 0, 1 \}}^\ZZ$ מידת מכפלה, $A \in \Ff$ מאורע אינווריאנטי, אם לכל $x$, $x \in A$ אם ורק אם $S(x) \in A$ כאשר $S(x)$ הזזה של הסדרה.
	אם $A$ אינווריאנטי אז $\PP(A) \in \{0, 1\}$.
\end{theorem}

\subsection{משתנים מקריים}
אם $(\Omega, \Ff, \PP)$ מרחב הסתברות אז משתנה מקרי הוא פונקציה מדידה מ־$\Omega$ ל־$\RR$.
תנאי מספיק למדידות הפונקציה $X$ הוא שלכל $t \in \RR$ יתקיים $X^{-1}((-\infty, t)) \in \Ff$a.
אם $X, Y$ משתנים מקריים אז $(X, Y)$ פונקציה מדידה מ־$\Omega$ ל־$\RR^2$, ואם $X_1, X_2, \ldots$ סדרת משתנים מקריים אז גם $(X_1, X_2, \ldots)$ פונקציה מדידה על $\RR^\NN$.

אם נחזור לדוגמה של קודם נוכל להגדיר את $(\Omega, \Ff, \PP)$ לפי הגרף $G = (V, E)$. נניח שלכל $e \in E$ יש משתנה מקרי $X_e$ המקריים $\PP(X^{-1}(\{0, 1\})) = 1$.
נקבל כך ש־${\{ x_e \}}_{e \in E}$ יוצר וקטור מקרי מ־$\Omega$ ל־${\{ 0, 1 \}}^E$.
\begin{corollary}
	כל פונקציה רציפה היא מדידה (בורל), והרכבה משמרת מדידות, לכן נסיק שאם $X, Y$ משתנים מקריים אז גם $X + Y, X \cdot Y$ משתנים מקריים, $\max(X, Y)$ וכן הלאה.
	אם ${\{ X_n \}}_{n \in \NN}$ סדרת משתנים מקריים אז גם $\sup_n X_n$ משתנה מקרי, שכן $\forall n,\ X_n \le t \iff \sup_n X_n \le t$.
	נסיק שגם $\limsup_{n \to \infty} X_n$ פונקציה מדידה ולכן משתנה מקרי גם כן.
	לכן גם $\{ w \in \Omega \mid \limsup_{n \to \infty} X_n = \liminf_{n \to \infty} X_n \}$ מדידה.
\end{corollary}
$X$ משתנה מקרי אז ההתפלגות של $X$ היא הדחיפה קדימה של $\PP$ דרך $X$, $\PP_X$.
נגדיר את התוחלת על־ידי,
\[
	\EE(X)
	= \int_{\Omega} X\ d\PP
	= \int_\RR x\ d\PP_X(x)
\]
ראינו גם מה קורה כשאנחנו מרכיבים פונקציה על $X$ במקרה זה. \\
אם $X, Y$ משתנים מקריים אז נאמר שהם בלתי־תלויים אם $\sigma(X), \sigma(Y)$ בלתי־תלויות, כאשר $\sigma(X) = \{ X^{-1}(B) \mid B \in \Ff \}$.
זה שקול ל־$\PP_{(X, Y)} = \PP_X \times \PP_Y$, ואם הם בלתי־תלויים אז גם $\EE(XY) = \EE(X) \EE(Y)$.

נאמר ש־$X_1, \ldots, X_n$ בלתי־תלויים בזוגות אם לכל $i \ne j$ $X_i, X_j$ בלתי־תלויים.
אם זה מתקיים אז גם $\var(\sum_{i = 1}^n X_i) = \sum_{i = 1}^n \var(X_i)$.
\begin{lemma}
	אם ${\{ A_n \}}$ בלתי־תלויים בזוגות ו־$\sum_{n = 1}^\infty \PP(A_n) = \infty$ אז $\PP(\limsup_{n \to \infty} A_n) = 1$.
\end{lemma}
\begin{proof}
	נגדיר $X_n = \indicator_{A_n}$ בלתי־תלויים בזוגות ו־$Y_n = \sum_{i = 1}^n X_i$.
	נראה שמתקיים,
	\[
		\PP({(\bigcup_{i = 1}^n A_i)}^C)
		= \PP(Y_n = 0)
		\to 0
	\]
	מתקיים מצ'בישב,
	\[
		\PP(Y_n = 0)
		\le \PP(| Y_n - \EE(Y_n) |)
		\le \frac{\var(Y_n)}{{(\EE(Y_n))}^2}
	\]
	וכן,
	\[
		\var(Y_n)
		= \sum_{i = 1}^n \var(Y_i)
		= \sum_{i = 1}^n p_i(1 - p_i)
		\le \sum_{i = 1}^n p_i
		= \EE(Y_n)
	\]
	ומשילוב שתי הטענות נקבל ש־$\PP(Y_n = 0) \to 0$, לכן $\PP({(\bigcup_{i = m}^\infty A_i)}^C) = 0$.
\end{proof}
\begin{exercise}
	נניח ש־$A_1, \ldots, A_n$ מאורעות בלתי־תלויים בזוגות כך ש־$\PP(A_i) = \frac{1}{2}$ לכל $i \le n$.
	מה יכולה להיות $\PP(\bigcap_{i = 1}^n A_i)$?
\end{exercise}

\section{שיעור 4 --- 6.5.2026}
\subsection{התכנסות}
בשיעור הקודם דיברנו על $\sigma$־אלגברות אי־תלות משתנים מקריים, בורל קנטלי השני בהנחת אי־תלות בזוגות.
עתה נדבר על התכנסות.
\begin{definition}[התכנסות כמעט תמיד]
	$X_n \xrightarrow{\text{a.s}} X$ אם $\PP(X_n \to X) = 1$.
\end{definition}
\begin{definition}[התכנסות בהסתברות]
	נסמן $X_n \xrightarrow{p} X$ אם $\lim_{n \to \infty} \PP(|X_n - X| > \varepsilon) = 0$ לכל $\varepsilon > 0$.
\end{definition}
במידה ראינו משפטי התכנסות על אינטגרלי לבג, ננסח תוצאות שלהם בניסוח הסתברותי.
\begin{theorem}[התכנסות תוחלת]
	$X_n \xrightarrow{\text{a.s.}} X$ אז,
	\begin{enumerate}
		\item קיים $M$ כך ש־$|X_n| \le M$ לכל $n$ (כמעט־תמיד) גורר ש־$\EE(X_n) \to \EE(X)$
		\item אם $0 \le X_1 \le \cdots$ אז $\EE(X_n) \to \EE(X)$ (במובן הרחב גם)
		\item אם קיים $Y$ עם תוחלת סופית כך ש־$|X_n| \le Y$ לכל $n$ אז $\EE(X_n) \to \EE(X)$
	\end{enumerate}
\end{theorem}
\begin{proposition}
	אם $X_n \to X$ כמעט תמיד אז $X_n \to X$ בהסתברות.
\end{proposition}
\begin{proof}
	נגדיר $A_n^\varepsilon = \{ |X_n - X| \le \varepsilon \}$, אז $X_n \xrightarrow{\text{a.s.}} X$ אם ורק אם,
	\[
		\forall \varepsilon > 0,\ \PP(\liminf_{n \to \infty} A_n^\varepsilon) = 1
	\]
	לעומת זאת, $X_n \xrightarrow{p} X$ אם ורק אם $\PP(A_n^\varepsilon) \to 1$ לכל $\varepsilon > 0$.
	מהלמה של פאטו נקבל ש־$1 = \PP(\liminf A_n^\varepsilon) \le \liminf \PP(A_n^\varepsilon) \le 1$
\end{proof}
הכיוון ההפוך של הגרירה הזאת הוא לא נכון, נראה דוגמה נגדית.
\begin{example}
	נגדיר $X_n = \indicator_{A_n}$ עבור $A_n$ בלתי־תלויים כך ש־$\PP(A_n) = p_n$ כך ש־$p_n \to 0$ ו־$\sum_{n = 1}^\infty p_n = \infty$ (לדוגמה $p_n = \frac{1}{n}$).
	אז $X_n \xrightarrow{p} 0$ אבל לא $X_n \xrightarrow{\text{a.s.}} 0$ מבורל קנטלי.
\end{example}
\begin{proposition}
	אם $X_n \xrightarrow{p} X$ אז קיימת תת־סדרה $X_{n_k}$ המקיימת $X_{n_k} \xrightarrow{\text{a.s.}} X$.
\end{proposition}
\begin{proof}
	נגדיר $B_n^\varepsilon = \{ |X_n - X| > \varepsilon \}$ ונסמן $p_n = \PP(B_n)$ אז $p_n \to 0$ לפי הגדרת התכנסות בהסתברות.
	נבחר תת־סדרה $n_k$ כך ש־$\sum_{k = 1}^\infty p_{n_k} < \infty$ ואז מבורל־קנטלי הראשון נקבל $|X_n - X| \le \varepsilon$ החל ממקום מסוים ולכן,
	\[
		X - \varepsilon \le \liminf X_{n_k} \le \limsup X_{n_k} \le X + \varepsilon
	\]
	זה נכון לכל $\varepsilon > 0$ ולכן יש שוויון, השאר נובע מטיעון אלכסון של סדרות.
\end{proof}
\begin{definition}[מטריקה על משתנים מקריים]
	עבור $X, Y$ משתנים מקריים מעל $(\Omega, \Ff, \PP)$ נגדיר את,
	\[
		d(X, Y) = \inf\{ \varepsilon > 0 \mid \PP(|X - Y| > \varepsilon) < \varepsilon \}
	\]
\end{definition}
\begin{exercise}
	הראו כי זוהי אכן מטריקה, והיא מגדירה $X_n \xrightarrow{p} X$.
\end{exercise}
אין מרחב מטרי שבו התכנסות סדרת משתנים מקריים שקולה להתכנסות כמעט תמיד.
ננסח טענה טופולוגית שלא נוכיח,
\begin{proposition}
	במרחב טופולוגי או מטרי $a_n \to a$ אם ורק אם לכל תת־סדרה $n_k$ יש תת־תת־סדרה $a_{n_{k_i}} \to a$.
\end{proposition}
אם נבחן את $X_n$ מהדוגמה אז $X_n \not\xrightarrow{\text{a.s.}} X$ אבל $X_n \xrightarrow{p} 0$ ולכן $X_{n_k} \xrightarrow{p} 0$ לכל תת־סדרה, ולפי הטענה שהוכחנו על תתי־סדרות מתכנסות,
נוכל לקחת תת־סדרה $X_{n_{k_i}} \xrightarrow{\text{a.s.}} 0$ ונקבל את המסקנה הבאה,
\begin{corollary}
	התכנסות כמעט תמיד איננה טופולוגית או מטרית.
\end{corollary}

\subsection{החוק החזק והחלש של המספרים הגדולים}
\begin{theorem}[החוק החזק של המספרים הגדולים]
	אם $X_n$ משתנים מקריים שווי־התפלגות ובלתי־תלויים ובעלי תוחלת $\mu$ אז,
	\[
		\frac{1}{n} S_n
		\xrightarrow{p} \mu
	\]
	כאשר $S_n = \sum_{i = 1}^n X_i$.
\end{theorem}
\begin{proof}[הוכחה (תחת אי־תלות בזוגות)]
	נניח שקיימת שונות למשתנים המקריים $\sigma^2$,
	\[
		\var(\frac{1}{n} \sum_{i = 1}^n X_i)
		= \frac{1}{n^2} \sum_{i = 1}^n \var(X_i) + \overbrace{\cov(\cdots)}^{= 0}
		= \frac{\sigma^2}{n}
		\to 0
	\]
	בהתאם,
	\[
		\PP(\left\lvert \frac{1}{n} \sum_{i = 1}^n X_i - \mu \right\rvert > \varepsilon)
		\le \frac{\var(\frac{1}{n} \sum_{i = 1}^n X_i)}{\varepsilon^2}
		= \frac{\sigma^2}{\varepsilon^2 n}
	\]
	מצ'בישב.

	עתה נניח רק ש־$X_i \ge 0$.
	יהי $M > 0$ כלשהו, נגדיר $X_i' = X_i \indicator_{\{ X_i \le M \}}$ ו־$X_i'' = X_i \indicator_{\{ X_i > M \}}$ ולכן $X_i = X_i' + X_i''$.
	נקבל,
	\[
		0 \le X_i' \le M
	\]
	נסמן $\EE(X_i') = \mu', \EE(X_i'') = \mu''$ ולכן גם $\mu = \mu' + \mu''$, ונבחין כי נקבל של־$X_i'$ שונות סופית ${(\sigma_i')}^2$. \\
	נראה שלכל $\varepsilon > 0$ קיים $M$ כך ש־$\EE(X_i'') \le \varepsilon$.
	נבחן את הסדרה $X_1 \indicator_{\{ X_1 \le M \}}$ אז היא מתכנסת מונוטונית ל־$X_i$ ולכן $\EE(X_1 \indicator_{\{ X_1 \le M \}}) \to \mu$ אז $\EE(X_1 \indicator_{\{ X_1 > M\}}) \to 0$. \\
	נסמן $S_n' = \sum_{i = 1}^n X_i'$ ואת $S_n''$ באופן דומה, ויהי $\varepsilon > 0$,
	\begin{align*}
		\PP(|\frac{1}{n} S_n - \mu| > \varepsilon)
		& = \PP(|S_n - n \mu| > \varepsilon n) \\
		& \le \PP(|S_n' - n \mu'| > \frac{\varepsilon}{2} n) + \PP(|S_n'' - n \mu''| > \frac{\varepsilon}{2} n) \\
		& \overset{(1)}{\le} \frac{{(\sigma')}^2}{{(\frac{\varepsilon}{2})}^2 n} + \frac{n \mu''}{\frac{\varepsilon}{2} n}
	\end{align*}
	כאשר $(1)$ נובע מצ'בישב ומרקוב.
	נבחר $M$ כך ש־$\mu'' \le {(\frac{\varepsilon}{2})}^2$
	ונקבל,
	\[
		\frac{{(\sigma')}^2}{{(\frac{\varepsilon}{2})}^2 n} + \frac{\mu''}{\frac{\varepsilon}{2}}
		< \frac{\varepsilon}{2} + \frac{\varepsilon}{2}
	\]

	בשלב האחרון נניח ש־$X$ בעל תוחלת ונפרק $X_i^+ = X_i \indicator_{\{ X_i \ge 0 \}}$ וכן $X_i^- = X_i \indicator_{\{ X_i \le 0 \}}$ ולכן $X_i = X_i^- + X_i^+$.
	בהתאם גם $\mu = \mu^- + \mu^+$.
	בהתאם גם $\frac{1}{n} \sum_{i = 1}^n X_i^+ \xrightarrow{p} \mu^+$ ואותו הדבר על $\mu^-$.
\end{proof}
עתה נעבור לחוק החזק.
\begin{theorem}[החוק החזק של המספרים הגדולים]
	אם $X_n$ משתנים מקריים שווי־התפלגות ובלתי־תלויים, אם הם בעלי תוחלת $\mu$ אז $\frac{1}{n} \sum_{i = 1}^n \xrightarrow{\text{a.s.}} \mu$.
\end{theorem}
\begin{proof}
	נניח תוחלת אפס בלי הגבלת הכלליות על־ידי מעבר למשתנה מקרי ממורכז.
	נבחן את $A_{n^2}^\varepsilon$, אז $\sum_{n = 1}^\infty \PP(A_{n^2}^\varepsilon) < \infty$ ולכן,
	\[
		\PP(\lim_{n \to \infty} A_{n^2}^\varepsilon) = 0
		\iff \PP(\liminf A_{n^2}^\varepsilon) = 1
	\]
	אם $n^2 \le k \le {(n + 1)}^2$ אז,
	\[
		|\sum_{i = n^2 + 1}^k X_i|
		= |S_k - S_{n_k^2}|
		\le M \cdot 2 n_k
	\]
	ולכן,
	\[
		\left\lvert \frac{S_k}{k} - \frac{S_{n_k^2}}{n_k^2} \cdot \frac{n_k^2}{k} \right\rvert
		\le \frac{M 2 n_k}{k}
		\to 0
	\]

	נראה שאם $X_i \ge 0$ שווי־התפלגות עם תוחלת $\mu$ רז $\PP(\limsup \frac{S_n}{n} \ge 4 \mu) = 0$.
	יהי $k \in \NN$ ונרצה לחסום את $\PP(\frac{S_{2^k}}{2^k} > 2 \mu)$.
	נגדיר,
	\[
		X_i' = X_i \indicator_{\{ X_i \le 2^k \}},
		\quad
		X_i'' = X_i \indicator_{\{ X_i > 2^k \}}
	\]
	אז,
	\[
		\PP(S_{2^k} > 2 \mu 2^k)
		\le \PP(S_{2^k}' > 2 \mu 2^k) + \PP(S_{2^k}'' > 0)
	\]
	נחסום אותם בנפרד,
	\begin{align*}
		\PP(S_{2^k}' > 2 \mu 2^k)
		& \le \PP(S_{2^k}' - 2^k \mu' > 2 \mu 2^k - 2^k \mu') \\
		& \le \PP(|S_{2^k}' - 2^k \mu'| \ge 2^k \mu) \\
		& \overset{\text{צ'בישב}}{\le} \frac{\var(S_{2^k}')}{2^{2^k} \mu^2} \\
		& = \frac{\var(X_i')}{2^k \mu^2} \\
		& \le \frac{\EE(X_1^2 \indicator_{\{ X_1 \le 2^k \}})}{2^k \mu^2}
	\end{align*}
	קיבלנו אם כך מאי־השוויון הקודם,
	\[
		\PP(S_{2^k} > 2 \mu 2^k)
		\le \PP(S_{2^k}' > 2 \mu 2^k) + \PP(S_{2^k}'' > 0)
		\le \frac{\EE(X_1^2 \indicator_{\{ X_1 \le 2^k \}})}{2^k \mu^2} + 2^k \PP(X_1 > 2^k)
	\]
	נסמן $p_k = \PP(2^k < X_1 \le 2^{k + 1})$ לכל $k \in \NN$ ועבור $k = 0$ נסמן $p_0 = \PP(X_1 \le 2)$.
	בהתאם,
	\[
		\sum_{k = 1}^\infty 2^k p_k
		\le 2 \mu
	\]
	מתקיים,
	\[
		\PP(X_1 > 2^k)
		= \sum_{j = k}^\infty p_j,
		\quad
		\EE(X_1^2 \indicator_{\{ X_1 \le 2^k \}})
		\le 4 \sum_{j = 0}^{k - 1} p_j 2^{2j} 
	\]
	וכן,
	\[
		\sum_{k = 1}^\infty \PP(X_1 > 2^k)
		= \sum_{k = 1}^\infty 2^k \sum_{j = k}^\infty p_j
		= \sum_{j = 1}^\infty p_j \sum_{k = 1}^j 2^k
		\le 2 \sum_{j = 1}^\infty p_j 2^j
		\le 4 \mu
		< \infty
	\]
	כלומר הטור סופי.
	\[
		\sum_{k = 1}^\infty \frac{1}{2^k} \sum_{j = 0}^{k - 1} p_j 2^{2j}
		= \sum_{j = 0}^\infty p_j 2^{2 j} \sum_{k = j}^\infty 2^{-k}
		= \sum_{j = 0}^\infty p_j 2^{2 j} \cdot 2^{-j} \cdot 2
		= \sum_{j = 0}^\infty p_j 2^j
		< \infty
	\]
	אז במסתכם מצאנו ש־$\sum_{k = 1}^\infty PP(\frac{1}{2^k} 2_{2^k} > 2 \mu) < \infty$, לכן מבורל־קנטלי הראשון נקבל שממקום מסוים,
	\[
		\frac{S_{2^k}}{2^k}
		\le 2 \mu
	\]
	נשים לב שאם $2^k < n \le 2^{k + 1}$ אז $S_{2^k} \le S_n \le S_{2^{k + 1}}$ ולכן,
	\[
		\frac{S_n}{n}
		\le \frac{S_{2^{k + 1}}}{n}
		= \frac{S_{2^{k + 1}}}{2^{k + 1}} \cdot \frac{2^{k + 1}}{n}
		\le 2 \cdot \frac{S_{2^{k + 1}}}{2^{k + 1}}
	\]
\end{proof}

\section{שיעור 5 --- 13.5.2026}
\subsection{התכנסות --- המשך}
נעבור להוכחת המשפט החזק לאחר שסיימנו את הלמה הרלוונטית בשיעור הקודם.
\begin{proof}
	נניח ש־$X_i \ge 0$ ויהי $\varepsilon > 0$, נבחר $M$ מספיק גדול כך שאם,
	\[
		X_i' = X_i \indicator_{\{ X_i \le M \}},
		\qquad
		X_i'' = X_i \indicator_{\{ X_i > M \}}
	\]
	אז 
	\[
		\EE(X_i') > \mu - \varepsilon,
		\qquad
		\EE(X_i'') < \varepsilon
	\]
	ונסמן,
	\[
		S_n = S_n' + S_n'',
		\qquad
		S_n' = \sum_{i = 1}^n X_i',
		\qquad
		S_n'' = \sum_{i = 1}^n X_i''
	\]
	ועתה נקבל,
	\begin{align*}
		\mu - \varepsilon
		& < \EE(X_1') \\
		& = \lim_{n \to \infty} \frac{S_n'}{n} \\
		& = \liminf_{n \to \infty} \frac{S_n'}{n} \\
		& \le \liminf_{n \to \infty} \frac{S_n}{n} \\
		& \le \limsup_{n \to \infty} \frac{S_n}{n} \\
		& \le \limsup_{n \to \infty} \frac{S_n'}{n} + \lim_{n \to \infty} \frac{S_n''}{n} \\
		& \le \EE(X_1') + 4 \EE(X_1'') \\
		& \le \mu + 4 \varepsilon
	\end{align*}
	עבור $X_i$ כלליים נפרק $X_i^+ = X_i \indicator_{\{X_i \ge 0\}}$ וכן $X_i^- = X_i \indicator_{\{ X_i < 0 \}}$ ואז $X_i = X_i^+ - X_i^-$.
\end{proof}
אם $\frac{S_n}{n} \xrightarrow{\text{a.s.}} a$ אז $a = \EE(X_i)$,
מהצד השני אם ל־$X_i$ אין תוחלת אז $\frac{S_n}{n}$ לא מתכנסת.
אם $0 \le X_i$ ו־$\EE(X_i) = \infty$ אז $\frac{S_n}{n} \to \infty$ כמעט תמיד.
קיימים משתנים מקריים $X_i$ בלתי־תלויים ושווי התפלגות ללא תוחלת כך ש־$\frac{S_n}{n} \xrightarrow{p} a \in \RR$.

\subsection{פרדוקס המעטפות}
נניח שאנו מקבלים שתי מעטפות, באחת יש $a$ כסף ובשנייה $2a$,
נקבל מעטפה אחת באקראי וראינו בה $b$, אז בשנייה יש או $2b$ או לחילופין $\frac{b}{2}$ כסף, ובהסתברות שווה. \\
בהתאם תוחלת המעטפה השנייה היא $\frac{1}{2} \cdot 2b + \frac{1}{2} \cdot \frac{b}{2} = \frac{5b}{4} > b$, כלומר לכאורה עלינו תמיד לבקש להחליף את המעטפה.

הפתרון הוא שאנחנו לא יכולים שלא להגדיר באופן ברור מהו $a$, גם לו יש התפלגות ואנו צריכים להתחשב בה.
הפעם נסמן $X \ge 0$ סכום הכסף במעטפה ובהתאם $2X$ במעטפה השנייה, כאשר $X$ משתנה מקרי, ונניח שהוא בדיד.
נסמן את המדידה של המעטפה הראשונה $Y$.
שתי האפשרויות בהתאם הן ש־$Y = X$ או ש־$Y = 2X$. \\
נניח ש־$\PP(X = 1) = p_a > 0$.
נרצה לבדוק אם יש התפלגות בדידה על $X$ כך שנקבל תוצאה נכונה.
אנו מניחים שמתקיים,
\[
	\PP(Y = a)
	= \PP(X = a) \frac{1}{2} + \PP(X = \frac{a}{2}) \frac{1}{2}
\]
ורוצים $\PP(X = a \mid Y = a) = \frac{1}{2}$.
נחשב,
\[
	\PP(X = a \mid Y = a)
	= \frac{\PP(X = a, Y = a)}{\PP(Y = a)}
	= \frac{\PP(X = a) \frac{1}{2}}{\PP(X = a) \frac{1}{2} + \PP(X = \frac{a}{2}) \frac{1}{2}}
	= \frac{p_a}{p_a + \frac{p_a}{2}}
	\implies p_{\frac{a}{2}} = p_a
\]
וזו כמובן סתירה.

עתה נניח ש־$Z \sim \operatorname{Geo}(\frac{1}{2})$, בהתאם,
\[
	\forall k \in \NN,\ 
	\PP(Z = k)
	= 2^{-k}
\]
ונניח ש־$X = 3^Z$ וננסה לחשב שוב את ההתפלגות של המעטפה השנייה.
אם $Y = 3^k$ אז או ש־$X = 3^k$ או ש־$X = 3^{k - 1}$.
הפעם נקבל תוך הנחת $k > 1$,
\[
	\PP(X = 3^k \mid Y = 3^k)
	= \frac{\PP(X = 3^k) \frac{1}{2}}{\PP(X = 3^k) \frac{1}{2} + \PP(X = 3^{k - 1}) \frac{1}{2}}
	= \frac{\PP(Z = k)}{\PP(Z = k) + \PP(Z = k - 1)}
	= \frac{2^{-k}}{2^{-k} + 2^{-(k - 1)}}
	= \frac{1}{1 + 2}
	= \frac{1}{3}
\]
כלומר,
\[
	\PP(X = 3^k \mid Y = 3^k)
	= \begin{cases}
		\frac{1}{3} & k > 1 \\
		1 & k = 1
	\end{cases}
\]
לכן אם $k > 1$ אז $\EE(\text{second envelope} \mid Y = 3^k) = \frac{1}{3} \cdot 3^{k + 1} + \frac{2}{3} \cdot 3^{k - 1} = 3^k \cdot \frac{11}{9}$.
כלומר שוב קיבלנו שמשתלם להחליף מעטפה, שוב באופן פרדוקסלי.
הסיבה לפרדוקס היא שהתוחלת עצמה היא לא סופית, ולכן התוצאה לא משמעותית כפי שהיא נראית.

\subsection{משפט הקירוב של ויירשטראס}
ראינו בקורסים קודמים שהפולינומים צפופים ברציפות, נעסוק בנושא זה מזווית הסתברותית.
\begin{definition}[פולינום ברנשטיין]
	אם $f \in C([0, 1])$ אז נגדיר,
	\[
		B_n(x)
		= \sum_{k = 0}^n f\left(\frac{k}{n}\right) \binom{n}{k} x^k {(1 - x)}^{n - k}
	\]
	נקרא פולינום ברנשטיין מסדר $n$ של $f$.
\end{definition}
\begin{remark}
	מתקיים $B_n(p) = \EE(f(\frac{S_n}{n}))$ עבור $S_n \sim \operatorname{Bin}(p, n)$.
\end{remark}
\begin{proposition}
	$\lVert B_n - f \rVert_\infty \xrightarrow{n \to \infty} 0$.
\end{proposition}
\begin{proof}
	$f$ רציפה ולכן חסומה ונסמן $M < \infty$ חסם כלשהו, $f$ גם רציפה במידה שווה.
	מרציפות במידה שווה לכל $\varepsilon > 0$,
	\[
		\sup_{|x - y| < \varepsilon} |f(x) - f(y)| \to 0
	\]
	וכן,
	\begin{align*}
		\forall p \in [0, 1],\ 
		\lVert B_n(p) - f(p) \rVert
		& = |\EE(f(\frac{S_n}{n}) - f(p))| \\
		& \le \EE(|f(\frac{S_n}{n}) - f(p)|) \\
		& \le \PP(|\frac{S_n}{n} - p| \ge \varepsilon) \cdot 2M + \PP(|\frac{S_n}{n} - p| < \varepsilon) \cdot \EE(| f(\frac{S_n}{n}) - f(p) |) \\
		& \le \frac{p(1 - p)}{n \varepsilon^2} \cdot 2M + 1 \cdot M(\varepsilon) \\
		& \xrightarrow[\varepsilon = \frac{1}{\sqrt{n}}]{n \to \infty} 0
	\end{align*}
\end{proof}

\subsection{הימורים}
נסמן $p > \frac{1}{2}$ ונהמר $a > 0$ בהסתברות $p$, נפסיד $a$ בהסתברות $1 - p$.
נקבל $Y_i \sim \operatorname{Ber}(p)$ בלתי־תלויים.
נגדיר את $X_n = X_{n - 1} + a_n(2Y_n - 1)$ וכן $X_0 = 1$ ולכן,
\[
	\EE(X_1)
	= \EE(X_0) + a(2p - 1),
	\quad
	E(X_2)
	= \EE(X_1) + a_2(2p - 1)
\]
מתקיים,
\[
	\frac{1}{n} \sum_{i = 1}^n \log X_i - \log X_{i - 1}
	= \frac{1}{n} \log X_n
	= \log \sqrt[n]{X_n}
\]
ולכן במקרה שלנו,
\[
	f_p(a) = p \log (1 + a) + (1 - p) \log (1 - a),
	\quad
	f_p'(a) = \frac{p}{1 + a}  (1 - p) \frac{1}{1 - a},
	\quad
	f_p'(a) = 0
	\iff p(1 - a) = (1 - p) (1 + a)
\]
ולכן $a = 2p - 1$, כלומר מצאנו שלא משתלם להשקיע או להמר בכל פעם את כל כספנו, אלא משתלם להשקיע מספר מסוים וחלקי של הכסף.

\subsection{אי־תלות בזוגות}
נניח ש־$X_i \sim U(\{-1, 1\})$ עבור $i \in [n]$ בלתי־תלויים בזוגות, ונבדוק את,
\[
	\PP(\sum_{i = 1}^n X_i = n)
	\le \PP(|\sum_{i = 1}^n \ge n)
	\le \frac{\var \sum_{i = 1}^n X_i}{n^2}
	= \frac{1}{n}
\]
נרצה להבין אם זהו החסם הטוב ביותר לערך זה.

\begin{proposition}
	נגדיר $Y_i \sim U(\{-1, 1\})$ בלתי־תלויים עבור $i \in [k]$ ונגדיר לכל $\emptyset \ne I \subseteq [k]$ את,
	\[
		X_I
		= \prod_{i \in I} Y_i
	\]
	אז $X_I \sim U(\{-1, 1\})$ והם בלתי־תלויים בזוגות.
\end{proposition}
\begin{proof}
	לחלק הראשון מספיק להראות ש־$\EE(X_I) = 0$, אבל,
	\[
		\EE(X_I)
		= \EE(\prod_{i \in I} Y_i)
		= \prod_{i \in I} \EE(Y_i)
		= \prod 0
		= 0
	\]

	נראה שהם בלתי־תלויים בזוגות, יהיו $I \ne J$ לא ריקות, אז,
	\[
		X_I \cdot X_J
		= X_{I \triangle J}
	\]
	ולכן,
	\[
		\EE(X_I X_J)
		= 0
	\]
	אם $Z_1, Z_2 \sim U(\{-1, 1\})$ הם בלתי־מתואמים ($\EE(Z_1 Z_2) = 0$) אז הם בלתי־תלויים (תרגיל).
\end{proof}
אתה נשאל מה ערך,
\[
	\PP(\forall \emptyset \ne I \subseteq [k]\ X_I = 1)
	= \PP(\forall i \in [k]\ Y_i = 1)
	= \frac{1}{2^k}
	= \frac{1}{n + 1}
\]
כלומר מצאנו חסם יותר טוב באופן זניח לשאלה המקורית שלנו.

\section{שיעור 6 --- 20.5.2026}
\subsection{שרשרות מרקוב}
נתחיל מתיאור של מצב.
נניח שצפרדע עומדת על פרח ויכולה לעבור לפרח אחר לחזור וכן הלאה.
ההסתברות לעבור מפרח א' לפרח ב' היא $q$ וההסתברות לעבור מפרח ב' לפרח א' היא $p$, ואנו רוצים להבין מה הסיכוי שתהיה בפרח מסוים לאורך זמן.
נקבל בהתאם שאם $X_n$ משתנה מקרי המתאר את מיקום הצפרדע באיטרציה ה־$n$ אז מתקיים,
\[
	\PP(X_{n + 1} = W \mid X_t = E, X_{n - 1} = x_{n - 1}, \ldots, X_0 = x_0) = p
	\quad
	\PP(X_{n + 1} = E \mid X_t = W, X_{n - 1} = x_{n - 1}, \ldots, X_0 = x_0) = q
\]
ניזכר שאם $A, B, C$ מאורעות כך ש־$A \cap B = \emptyset$ ו־$\PP(C \mid A) = \PP(C \mid B) = p$ אז $\PP(C \mid A \cup B) = p$.
נקבל אם כך שמתקיים,
\[
	\PP(X_{n + 1} = W \mid X_n = E) = p,
	\qquad
	\PP(X_{n + 1} = E \mid X_n = W) = q
\]
בדיוק, כלומר אנו לא תלויים במה שקרה לפני האיטרציה האחרונה.
נסמן $\mu_n(E) = \PP(X_n = E), \mu_n(W) = \PP(X_n = W)$ ולכן,
\begin{align*}
	\mu_{n + 1}(E)
	& = \PP(X_{n + 1} = E) \\
	& = \PP(X_{n + 1} = E \mid X_n = W) \PP(X_n = W) + \PP(X_{n + 1} = E \mid X_n = E) \PP(X_n = E) \\
	& = q \mu_n(W) + (1 - p) \mu_n(E)
\end{align*}
ולכן גם $\mu_{n + 1}(W) = (1 - q) \mu_n(W) + p \mu_n(E)$.
אם נגדיר,
\[
	P = \begin{pmatrix}
		1 - p & p \\
		q & 1 - q
	\end{pmatrix}
\]
נקבל שמתקיים $\mu_{n + 1} = \mu_n \cdot P$ ולכן גם $\mu_n = \mu_0 P^n$ לכל $n \in \NN$.
אם נקבל התכנסות על $n \to \infty$ אז קיים $\mu = (r, 1 - r)$ כך שמתקיים $\mu = \mu P$ ולכן נובע ש־$r = (1 - p) r + q (1 - r)$ ולכן $r = \frac{q}{p + q}$.
בהתאם $\mu = (\frac{q}{p + q}, \frac{p}{p + q})$.
עתה נניח ש־$\mu_n \to \mu$ ולכן,
\[
	\Delta_{n + 1}
	= \mu_{n + 1}(E) - \frac{q}{p + q}
	= (1 - p) \mu_n(E) + q (1 - \mu_n(E)) - \frac{q}{p + q}
	= q + (1 - p - q) (\mu_n(E) - \frac{q}{p + q})
\]
ונובע,
\[
	\Delta_{n + 1} = (1 - p - q) \Delta_n
\]
ולכן גם,
\[
	\Delta_n = (1 - p - q) \Delta_0
\]
$X_0, X_1, \ldots$ המשתנים המקריים שמקבלים ערכים ב־$S$ קבוצת המצבים עבור $S$ בדידה עם טופולוגיה ־$\sigma$־אלגברה בדידה.
\begin{definition}[שרשרת מרקוב]
	$X_0, X_1, \ldots$ היא שרשרת מרקוב עם מטריצת מעבר $P \in M_{S \times S}$ אם,
	\[
		\PP(X_{t + 1} = s_{t + 1} \mid X_t = s_t, X_{t - 1} = s_{t - 1}, \ldots, X_0 = s_0) = P_{s_{t + 1}, s_t}
	\]
	לכל $t$ ולכל $s_0, \ldots, s_{t + 1} \in S$.
\end{definition}
למעשה,
\[
	P(s_{t + 1}, s_t)
	= \PP(X_{t + 1} = s_{t + 1} \mid X_t = s_t)
\]
בלבד מאותה הסיבה של קודם.
התפלגות על $S$ היא וקטור עם ערכים אי־שליליים שסכומם $1$, ו־$\Delta_S$ נקרא סימפלקס קבוצת כל ההתפלגויות על $S$.
\begin{definition}[סטציונרית]
	$\mu \in \Delta_S$ נקראת סטציונרית אם $\mu = \mu P$.
\end{definition}
\begin{proposition}
	אם $S$ סופית אז קיימת התפלגות סטציונרית.
\end{proposition}
\begin{proof}
	ניקח $\mu_0 \in \Delta_S$ כלשהי ונגדיר $\mu_t = \mu_0 P^t$, נגדיר,
	\[
		\Delta_S \ni
		\nu_t
		= \frac{1}{t} \sum_{i = 1}^t \mu_i
	\]
	אבל $\Delta_S$ קומפקטית ולכן יש תת־סדרה מתכנסת $\nu_{t_k} \to \nu$, נרצה להראות ש־$\nu$ היא סטציונרית.
	\[
		\nu_t P
		= \frac{1}{t} \sum_{i = 1}^y \mu_i P
		= \frac{1}{t} \sum_{i = 2}^{t + 1} \mu_i
	\]
	ונקבל,
	\[
		\nu(P - I)
		= \nu_t P - \nu_t
		= \frac{1}{t} (\nu_{t + 1} - \nu_t)
	\]
	ולכן $\nu(P - I) \to 0$ ובפרט $\nu_{t_k} (P - I) \to 0$ ולכן $\nu(P - I) = \nu P - \nu$.
\end{proof}
\begin{example}
	יהי $G = (V, E)$ גרף פשוט לא מכוון עם דרגות סופיות.
	מטריצת המעבר של הילוך מקרי פשוט על $G$ היא,
	\[
		P(u, v)
		= \begin{cases}
			\frac{1}{\deg u} & u \sim v \\
			0 & \text{otherwise}
		\end{cases}
	\]
\end{example}
\begin{proposition}
	המידה $\mu(v) = \deg v$ היא סטציונרית, $\mu = \mu P$.
\end{proposition}
\begin{proof}
	\[
		(\mu P)(v)
		= \sum_{u \in V} \mu(u) P(u, v)
		= \sum_{u \sim v} \deg(u) \cdot \frac{1}{\deg u}
		= \deg v
		= \mu(v)
	\]
\end{proof}
אם $G$ סופי אז אפשר לנרמל ולקבל התפלגות ב־$\Delta_V$, אם $G$ אינסופי אז זה לא עובד. \\
נרחיב עתה את הדוגמה ונניח ש־$G$ בעל משקולות, כלומר קיימת $C : E \to \RR^+$ ונגדיר,
\[
	\pi(u)
	= \sum_{u \sim v} C_{uv}
\]
ונגדיר,
\[
	P(u, v)
	= \frac{C_{uv}}{\pi(u)}
\]
\begin{proposition}
	זוהי מטריצה מקרית ו־$\pi$ היא סטציונרית, כלומר $\pi = \pi P$.
\end{proposition}

\section{שיעור 7 --- 27.5.2026}
\subsection{שרשרות מרקוב}
נמשיך את הדיון על שרשרות בהתאם להגדרה שסיפקנו בהרצאה הקודמת.
נניח ש־$X_0, \ldots$ סדרת משתנים מקריים שמקבלים ערכים בקבוצה $S$ הבדידה.
הגדרנו גם את $P \in M_{S \times S}$ כך שלכל $s_0, \ldots, s_1 \in S$ מתקיים,
\[
	\PP(X_{n + 1} = s_{n + 1} \mid \forall i \le n,\ X_i = s_i)
	= P(s_n, s_{n + 1})
\]
ולכן בפרט גם המטריצה היא סטוכסטית, כלומר $\forall x,\ \sum_{y \in S} P(x, y) = 1$.
גם הראינו שאם $\mu_0 \in \Delta_S$ ו־$\mu_0 \sim X_0$ אז $X_n \sum \mu_n = \mu_0 P^n$.
הוכחנו שלכל $P$ סטוכסטית קיימת $\pi \in \Delta_S$ כך ש־$\pi = \pi P$ (סטציונרית) על $S$ סופית. \\
עתה נרצה להראות שמתקיים,
\[
	\PP_s(X_n = t)
	= P^n(s, t)
\]
\[
	\PP_s(X_{n + 1} = t)
	= \sum_{r \in S} \PP_s(X_n = n) \PP_s(X_{n + 1} t \mid X_n = r)
	= \sum_{r \in S} P^n(s, r) P(r, t)
	= (P^n \cdot P)(s, t)
	= P^{n + 1}(s, t)
\]
\begin{definition}[איברים ניתנים להגעה]
	בהינתן $P$ נאמר שאפשר להגיע מ־$x$ ל־$y$ עבור $x, y \in S$ אם קיים $n$ כך ש־$P^n(x, y) > 0$.
\end{definition}
\begin{definition}[שרשרת אי־פריקה]
	שרשרת נקראת אי־פריקה אם אפשר להגיע מכל מצב לכל מצב
\end{definition}
\begin{theorem}[יחידות התפלגות אי־פריקה]
	אם $P$ אי־פריקה ו־$S$ סופית אז ההתפלגות הסטציונרית היא יחידה.
\end{theorem}
\begin{proof}
	נראה שאם $P$ אי־פריקה אז מימד המרחב העצמי של $1$ הוא $1$.
	נסמן את $h = (1, \ldots, 1)$ וקטור עצמי כזה, כלומר $h = P h$.
	בהתאם,
	\[
		h(x)
		= \sum_y P(x, y) h(y)
		= \EE_x(h(X_1))
	\]
	נגדיר $M = \max_{x \in S} h(x)$ ונסמן $x_0$ כך ש־$h(x_0) = M$, אז,
	\[
		M
		= h(x_0)
		= \sum_u P(x_0, y) h(y)
		\le \sum_y P(x_0, y) M
		= M
	\]
	ולכן $P(x_0, y) = P(x_0, y) M$ לכל $y$.
	נובע שלכל $y$ עבורו $0 < P(x_0, y)$ מתקיים $P(y) = M$.
	$P$ אי־פריקה ולכן אפשר להגיע מכל $x$ לכל $y$ ונובע ש־$h \equiv M$.
\end{proof}
אנו יודעים ש־$\pi = \pi P$ ולכן ממשפט זה נסיק שהוא היחיד במקרה זה. \\
תהי $P$ שרשרת־מרקוב על $S$ סופית, עבור $x$ נסמן,
\[
	\NN
	\supseteq A_x
	= \{ n > 0 \mid P^n(x, x) < 0 \}
.\]
ונסמן $a_x = \gcd A_x$ להיות המחזור של $x$. \\
$A_x$ סגורה לחיבור כי $P^n P^m = P^{n + m}$ ולכן אם $0 < P^n(x, x)$ וגם $0 < P^m(x, x)$ אז גם $0 < P^{n + m}(x, x)$.
\begin{proposition}
	אם $A \subseteq \NN$ סגורה לחיבור אז $|\gcd(A) \NN \setminus A| < \infty$.
\end{proposition}
\begin{proposition}
	אם $P$ אי־פריקה אז כל ה־$a_x$ שווים.
\end{proposition}
\begin{proof}
	יהיו $x, y \in S$ ונרצה להראות ש־$a_x = a_y$. \\
	יש $n$ כך ש־$P^n(x, y) > 0$ ויש $m$ כך ש־$P^m(y, x) > 0$.
	אן $k \in A_x$ אז $m + n + k \in A_y$ כי,
	\[
		P^{n + m + k}(y, y) = P^m(y, x) P^k(x, x) P^n(x, y) > 0
	\]
	לכן $\gcd(A_y) \le \gcd(A_x)$ ומסימטריה $a_y = a_x$.
\end{proof}
\begin{corollary}
	אם $P$ אי־פריקה אז $\forall x,\ a_x = a$ לאיזשהו $a \in \NN$, נקרא לו המחזור של $P$.
\end{corollary}
\begin{theorem}
	אם $P$ אי־פריקה על $S$ סופית ולא מחזורית (כלומר $a = 1$) אז לכל $\mu_0 \in \Delta_S$ מתקיים $\mu_n = \mu_0 P^n \to \pi$ עבור $\pi$ המידה הסטציונרית היחידה.
\end{theorem}
\begin{proof}
	נניח ש־$P(x, y) > 0$ לכל $x, y$.
	במקרה זה קיים $\varepsilon > 0$ כך שלכל $x, y \in S$ מתקיים $P(x, y) \ge \varepsilon \pi(y)$.
	לכן אם $\mu_0 \in \Delta_S$ ונסמן $\mu_0 = \nu_0$ אז קיים $\nu_1 \in \Delta_S$ כך שמתקיים,
	\[
		\nu_0 P
		= \varepsilon \pi + (1 - \varepsilon) \nu_1
	.\]
	ולכן קיים גם $\nu_2$ כך שגם,
	\[
		\nu_1 P
		= \mu_1 P
		= \varepsilon \pi + (1 - \varepsilon) \nu_2
	.\]
	ונמשיך ככה לקבל סדרה,
	\[
		\nu_n P
		= \varepsilon \pi + (1 - \varepsilon) \nu_{n + 1}
	.\]
	נקבל,
	\[
		\mu_n
		= (1 - {(1 - \varepsilon)}^n) \pi + {(1 - \varepsilon)}^n \nu_n
	.\]
	ולכן,
	\begin{align*}
		\mu_{n + 1}
		& = \mu P \\
		& = (1 - {(1 - \varepsilon)}^n) \pi + {(1 - \varepsilon)}^n \nu_n P \\
		& = (1 - {(1 - \varepsilon)}^n) \pi + {(1 - \varepsilon)}^n (\varepsilon \pi + (1 - \varepsilon) \nu_{n + 1}) \\
		& = (1 - {(1 - \varepsilon)}^n - {(1 - \varepsilon)}^n \varepsilon) \pi + {(1 - \varepsilon)}^{n + 1} \nu_{n + 1}
	.\end{align*}

	עתה נניח ש־$P$ לא מחזורית ולכן קיים $n \in \NN$ כך ש־$\forall x, y \in S,\ P^n(x, y) > 0$ כי מאי־פריקות קיים $n$ כך ש־$0 < P^n(x, y)$ ומהטענה הקודמת מתקיים $0 < P^n(x, x)$ עבור כל $n$ גדול מספיק.
	לכן,
	\[
		P^{n + m}(x, y)
		\ge P^m(x, x) P^n(xm y) > 0
	.\]
	ונובע ש־$P^r(x, y) > 0$ לכל $x, y$ לאיזשהו $r \in \NN$.
	מתקיים $\pi = \pi P^r$ ולפי המקרה הקודם לשרשרת $P^r$ נקבל $\mu P^{r n} \to \pi$ ובאופן דומה גם $\mu P^{rn + k} = (\mu P^k) P^{r n} \to \pi$ לכל $0 \le k < r$.
\end{proof}
\begin{definition}[גרף קיילי]
	אם $G$ חבורה ו־$B \subseteq G$ סימטרית, כלומר $b \in B \implies b^{-1} \in B$.
	נגדיר גרף קיילי עם קבוצת יוצרים $B$ עם הקשר $x \sim y \iff x y^{-1} \in B$.
	אם $G$ סופית אז התפלגות אחידה היא סטציונרית.
\end{definition}

\listoftheorems[title=הגדרות ומשפטים,ignoreall,show={theorem,definition},swapnumber,onlynamed={proposition}]
\end{document}
